\begin{exercise}[Enthalpie libre d'un gaz de Joule]

$n$ mol de gaz, d’équation d’état :
\begin{align*}
P (V-n \cdot b) = n RT
\end{align*}
, suivent une isotherme de $(T_0, P_i)$ à $(T_0, P_f)$.

\begin{questions}
\question Écrire l’expression différentielle $dG$ de l’enthalpie libre $G(T,P)$.

\question Déterminer la variation d’enthalpie libre $\Delta G$ au cours de cette transformation.

\question Bien qu’il obéisse à la $1^e$ loi de Joule, ce gaz (dit \og de Joule \fg{}) ne vérifie pas la $2^e$ loi de Joule :
\begin{align*}
dH = n C_{P,m} dT + n b dP
\end{align*}

Vérifie-t-il la relation de Mayer du gaz parfait $C_{P,m} - C_{V,m} = R$ ?
\end{questions}

\end{exercise}